\section{Experiment}
The work presented here in this thesis is based on a series of measurements carried out by the \SI{5}{MeV} accelerator at Aarhus University in 2017. From these, specifically the measurements involving the \SI{17.76}{MeV} $0^+$ resonance of \ce{^12C} is of interest. Although this section is written from the perspective of the $0^+$ resonance, with some minor changes in the numbers, it is equally valid for any of the other resonances. 

\subsection{Experimental setup}
A standard \SI{5}{MeV} Van de Graaf accelerator is used to generate a \SI{2001}{keV} proton beam. The beam is then guided through a series of slits before impacting the target nucleus in a reaction chamber. A Faraday cup is placed further downstream to pick up any remnant beam particles. The target is a $\sim$\SI{120}{nm} thick \ce{^11B} layer on carbon backing. \\

As already mentioned in the introduction to this section, this setup was also used to measure other nearby excited states of carbon, such as the \SI{16.57}{MeV} $2^-$ and \SI{18.35}{MeV} $3^-$ resonances. These two sets of measurements will be used mainly to provide some perspective on the results of the $0^+$ resonance.

\subsubsection{Detectors}
Four double-sided silicon-strip detectors (DSSD) from \textit{Micron Semiconductor Ltd} were used; two each of the models W1 and S3. The front side of each detector is p-doped, and is split into a number of strips. Similarly the back side is n-doped, and split into strips running perpendicularly to the front side. Thus these two sets of strips effectively constitutes a grid of pixels, which allows us to sensitively measure the direction of a particle. \\

When a particle hits a detector, its surface atoms are ionized, and produces electron-hole pairs proportional to the energy of the incident particle. The bias voltage collects the electrons and holes on a pair of aluminum contacts, from which a signal can be measured. Since any energy deposited in the aluminum is lost, it effectively acts as a dead layer for the particles. \\

The \SI[product-units=power]{49.5x49.5}{mm} surface area of the W1 detectors are split into $16$ \SI[product-units=power]{3x50}{mm} rectangular strips on each side, thus creating $256$ \SI{9}{mm^2} pixels. They are \SI{60}{\mu m} thick, which gives them a dead layer equivalent to \SI{100(10)}{nm} of silicon \note{is this true?}. They are placed parallel to the beam axis, at a distance of \SI{4}{cm}. Due to the symmetry of the setup, it does not matter which names we associate with the left and right side, and they have thus been named simply Det1 and Det2. \\

The S3 detectors are annular, such that the beam can travel through the central hole. They have an inner radius of \SI{11}{mm}, and an outer radius of \SI{35}{mm}. The front side is split into 32 radial spokes, while the back side is split into 24 rings, thus constituting $768$ pixels of varying sizes. They are \SI{322}{\mu m} thick, which gives them a dead layer equivalent to \SI{570(57)}{nm} of silicon \note{is this true?}. The two detectors are both placed perpendicularly with the beam axis, at a separation of \SI{36}{mm} from the target. The upstream detector is named SU(pstream), while the downstream detector is named SD(ownstream). Due to its downstream position, the inner $16$ rings of the SD are covered by a collimator. Furthermore its spokes 1, 2, 3, and 32 have not registered anything, and are most likely broken. All other strips in the other detectors appear to be working correctly. \\

As can be seen on figure \note{make figure}, the setup covers a rather large solid angle, which is ideal for measuring multi-coincidence data. This will later allow us to filter the data to only the events where all three decay products were detected. To reduce the amount of data gathered, a trigger with a logical AND condition is used, thus requiring a simultaneous hit in at least two detectors. Furthermore, to reduce the number of detected Rutherford scattered protons, a downscaling is applied to each individual detector: For every four events accepted by the SU, two are accepted in each of the W3's, and only a single event is accepted by the SD due to its downstream position. \\

The analogue signals from the trigger logic are digitized by the data acquisition system. \note{more, mention S3 only BT no FT}